\documentclass[12pt,a4paper]{article}

% ============================================
% PAQUETES
% ============================================
\usepackage[utf8]{inputenc}
\usepackage[spanish]{babel}
\usepackage[T1]{fontenc}
\usepackage{amsmath,amsfonts,amssymb}
\usepackage{graphicx}
\usepackage{booktabs}
\usepackage{array}
\usepackage{multirow}
\usepackage{xcolor}
\usepackage{colortbl}
\usepackage{float}
\usepackage{subcaption}
\usepackage[colorlinks=true,linkcolor=blue,citecolor=blue,urlcolor=blue]{hyperref}
\usepackage{geometry}
\usepackage{fancyhdr}
\usepackage{listings}
\usepackage{tikz}
\usetikzlibrary{shapes,arrows,positioning,calc}
\usepackage{algorithm}
\usepackage{algorithmic}
\usepackage{tcolorbox}

\geometry{margin=2.5cm}

% Colores personalizados
\definecolor{loyolablue}{RGB}{0,102,153}
\definecolor{codegreen}{rgb}{0,0.6,0}
\definecolor{codegray}{rgb}{0.5,0.5,0.5}
\definecolor{codepurple}{rgb}{0.58,0,0.82}
\definecolor{backcolour}{rgb}{0.95,0.95,0.92}
\definecolor{melcolor}{RGB}{220,53,69}
\definecolor{goodgreen}{RGB}{40,167,69}

% Configuración de listings para código Python
\lstdefinestyle{pythonstyle}{
    backgroundcolor=\color{backcolour},
    commentstyle=\color{codegreen},
    keywordstyle=\color{codepurple},
    numberstyle=\tiny\color{codegray},
    stringstyle=\color{codegreen},
    basicstyle=\ttfamily\footnotesize,
    breakatwhitespace=false,
    breaklines=true,
    captionpos=b,
    keepspaces=true,
    numbers=left,
    numbersep=5pt,
    showspaces=false,
    showstringspaces=false,
    showtabs=false,
    tabsize=2,
    language=Python
}
\lstset{style=pythonstyle}

% Cajas de información
\newtcolorbox{conceptbox}[1][]{
    colback=blue!5!white,
    colframe=loyolablue,
    fonttitle=\bfseries,
    title=#1
}

\newtcolorbox{warningbox}[1][]{
    colback=red!5!white,
    colframe=melcolor,
    fonttitle=\bfseries,
    title=#1
}

\newtcolorbox{tipbox}[1][]{
    colback=green!5!white,
    colframe=goodgreen,
    fonttitle=\bfseries,
    title=#1
}

% Encabezado y pie de página
\pagestyle{fancy}
\fancyhf{}
\fancyhead[L]{\small Deep Learning - Clasificación de Imágenes}
\fancyhead[R]{\small Universidad Loyola}
\fancyfoot[C]{\thepage}
\renewcommand{\headrulewidth}{0.4pt}

\begin{document}

% ============================================
% PORTADA
% ============================================
\begin{titlepage}
    \centering
    \vspace*{1cm}

    {\Large \textbf{Universidad Loyola}}\\[0.5cm]
    {\large Máster Universitario en Inteligencia Artificial}\\[0.3cm]
    {\large Curso Académico 2025/2026}\\[2cm]

    \rule{\textwidth}{1.5pt}\\[0.5cm]
    {\Huge \textbf{Clasificación de Lesiones Cutáneas mediante Deep Learning}}\\[0.3cm]
    {\Large Bloque A: Modelos Unimodales con Imágenes}\\[0.5cm]
    \rule{\textwidth}{1.5pt}\\[2cm]

    {\large \textbf{Proyecto Final Integrado}}\\[0.3cm]
    {\large Machine Learning - Deep Learning}\\[2cm]

    \begin{tabular}{rl}
        \textbf{Autor:} & David More \\[0.3cm]
        \textbf{Profesores:} & Diego Canales Aguilera \\
        & Antonio M. Durán Rosal \\
    \end{tabular}\\[3cm]

    \vfill
    {\large Enero 2026}
\end{titlepage}

% ============================================
% ÍNDICE
% ============================================
\tableofcontents
\newpage

% ============================================
% 1. INTRODUCCIÓN
% ============================================
\section{Introducción}

\subsection{Contexto del Problema}

El cáncer de piel es uno de los tipos de cáncer más comunes a nivel mundial. Según la Organización Mundial de la Salud, cada año se diagnostican más de 3 millones de casos de cáncer de piel no melanoma y aproximadamente 132.000 casos de melanoma maligno. La detección temprana es \textbf{crítica}: un melanoma detectado en etapa temprana tiene una tasa de supervivencia superior al 95\%, mientras que en etapas avanzadas esta tasa cae drásticamente por debajo del 20\%.

\begin{warningbox}[Importancia Clínica del Melanoma]
El melanoma (MEL) representa solo el 1\% de los cánceres de piel, pero es responsable de la gran mayoría de muertes por cáncer de piel. Por esta razón, en nuestro trabajo \textbf{priorizamos la detección correcta de melanomas} sobre otras métricas generales.
\end{warningbox}

La dermatoscopia es una técnica no invasiva que permite visualizar estructuras de la piel no visibles a simple vista. Sin embargo, la interpretación de estas imágenes requiere años de experiencia especializada, y existe una variabilidad significativa entre dermatólogos. Aquí es donde el \textit{Deep Learning} puede aportar valor: sistemas de apoyo al diagnóstico que ayuden a los profesionales médicos.

\subsection{Dataset ISIC: Características}

El dataset utilizado proviene de la \textit{International Skin Imaging Collaboration} (ISIC), específicamente adaptado para clasificación multiclase. Nuestro subconjunto contiene \textbf{6 tipos de lesiones cutáneas}:

\begin{table}[H]
\centering
\caption{Clases del dataset y su significado clínico}
\label{tab:clases}
\begin{tabular}{@{}llp{7cm}@{}}
\toprule
\textbf{Código} & \textbf{Nombre Completo} & \textbf{Descripción} \\
\midrule
ACK & Queratosis Actínica & Lesión precancerosa causada por exposición solar \\
BCC & Carcinoma Basocelular & Cáncer de piel común, raramente metastásico \\
\rowcolor{red!10} MEL & \textbf{Melanoma} & \textbf{El más peligroso. Prioridad máxima de detección} \\
NEV & Nevus (Lunar) & Lesión benigna común \\
SCC & Carcinoma Espinocelular & Cáncer de piel con potencial metastásico \\
SEK & Queratosis Seborreica & Lesión benigna común en adultos mayores \\
\bottomrule
\end{tabular}
\end{table}

\subsection{Distribución del Dataset}

\begin{table}[H]
\centering
\caption{División del dataset por conjuntos}
\label{tab:split}
\begin{tabular}{@{}lccc@{}}
\toprule
\textbf{Conjunto} & \textbf{Imágenes} & \textbf{Proporción} & \textbf{Uso} \\
\midrule
Entrenamiento (Train) & 500 & 70\% & Aprendizaje de los modelos \\
Validación (Val) & 130 & 15\% & Ajuste de hiperparámetros \\
Test & 161 & 15\% & Evaluación final (nunca visto) \\
\bottomrule
\end{tabular}
\end{table}

\begin{warningbox}[Problema Crítico: Desbalanceo de Clases]
La clase MEL (Melanoma) tiene significativamente \textbf{menos muestras} que las demás clases ($\approx$35 imágenes de entrenamiento). Este desbalanceo es un desafío técnico importante que abordaremos con técnicas específicas.
\end{warningbox}

\subsection{Objetivos del Trabajo}

Los objetivos específicos de este bloque son:

\begin{enumerate}
    \item \textbf{Establecer una línea base} con arquitecturas CNN preentrenadas
    \item \textbf{Aplicar Transfer Learning} para compensar el tamaño limitado del dataset
    \item \textbf{Implementar técnicas de regularización} para evitar overfitting
    \item \textbf{Construir un ensemble} que combine múltiples arquitecturas
    \item \textbf{Priorizar la detección de melanoma} en las decisiones de diseño
\end{enumerate}

% ============================================
% 2. FUNDAMENTOS TEÓRICOS
% ============================================
\section{Fundamentos Teóricos}

Esta sección explica los conceptos fundamentales necesarios para entender las decisiones técnicas tomadas en el proyecto.

\subsection{Redes Neuronales Convolucionales (CNNs)}

Las \textbf{Redes Neuronales Convolucionales} son arquitecturas de Deep Learning especialmente diseñadas para procesar datos con estructura de cuadrícula, como las imágenes.

\subsubsection{¿Por qué Convoluciones para Imágenes?}

Una imagen de 224×224 píxeles en color tiene $224 \times 224 \times 3 = 150.528$ valores de entrada. Si conectáramos cada píxel a cada neurona de una capa densa de 1000 neuronas, tendríamos más de 150 millones de parámetros solo en esa capa. Esto es:

\begin{itemize}
    \item \textbf{Computacionalmente prohibitivo}
    \item \textbf{Propenso a overfitting} (demasiados parámetros para los datos disponibles)
    \item \textbf{Incapaz de capturar patrones locales} como bordes o texturas
\end{itemize}

Las convoluciones resuelven estos problemas mediante dos principios fundamentales:

\begin{conceptbox}[Principios de las Convoluciones]
\begin{enumerate}
    \item \textbf{Conectividad Local}: Cada neurona solo ``ve'' una pequeña región de la imagen (campo receptivo), no toda la imagen.
    \item \textbf{Compartición de Pesos}: El mismo filtro (kernel) se aplica a toda la imagen, reduciendo drásticamente los parámetros.
\end{enumerate}
\end{conceptbox}

\subsubsection{Operación de Convolución}

Matemáticamente, la convolución 2D se define como:

\begin{equation}
(I * K)[i,j] = \sum_{m}\sum_{n} I[i+m, j+n] \cdot K[m,n]
\end{equation}

Donde $I$ es la imagen de entrada, $K$ es el kernel (filtro) y $*$ denota la operación de convolución. El kernel ``desliza'' sobre la imagen, multiplicando elemento a elemento y sumando los resultados.

\begin{figure}[H]
\centering
\begin{tikzpicture}[scale=0.6]
    % Input matrix
    \node at (-1, 3) {\textbf{Entrada}};
    \foreach \i in {0,...,4} {
        \foreach \j in {0,...,4} {
            \draw (\i, -\j) rectangle (\i+1, -\j+1);
        }
    }
    % Highlight receptive field
    \fill[blue!30] (1,0) rectangle (4,-3);

    % Kernel
    \node at (7, 3) {\textbf{Kernel 3×3}};
    \foreach \i in {0,...,2} {
        \foreach \j in {0,...,2} {
            \draw (6+\i, -\j) rectangle (7+\i, -\j+1);
        }
    }
    \fill[blue!30] (6,0) rectangle (9,-3);

    % Arrow
    \draw[->, thick] (10, -1) -- (12, -1);

    % Output
    \node at (14.5, 3) {\textbf{Salida}};
    \foreach \i in {0,...,2} {
        \foreach \j in {0,...,2} {
            \draw (13+\i, -\j) rectangle (14+\i, -\j+1);
        }
    }
    \fill[green!30] (14,-1) rectangle (15,-2);

\end{tikzpicture}
\caption{Operación de convolución: el kernel 3×3 desliza sobre la entrada, produciendo un mapa de características de menor dimensión.}
\end{figure}

\subsubsection{Componentes de una CNN}

Una CNN típica contiene los siguientes tipos de capas:

\begin{enumerate}
    \item \textbf{Capas Convolucionales}: Extraen características mediante filtros aprendibles
    \item \textbf{Funciones de Activación (ReLU)}: Introducen no-linealidad: $f(x) = \max(0, x)$
    \item \textbf{Capas de Pooling}: Reducen la dimensionalidad espacial (ej. Max Pooling)
    \item \textbf{Batch Normalization}: Normaliza las activaciones para estabilizar el entrenamiento
    \item \textbf{Capas Fully Connected}: Al final, para la clasificación
\end{enumerate}

\subsection{ResNet: Redes Residuales}

\subsubsection{El Problema de la Degradación}

Antes de ResNet, entrenar redes muy profundas era problemático. Contraintuitivamente, añadir más capas a una red \textbf{empeoraba} el rendimiento, incluso en el conjunto de entrenamiento. Esto no era overfitting, sino un problema de optimización conocido como \textbf{degradación}.

\subsubsection{La Solución: Conexiones Residuales}

Kaiming He et al. (2015) propusieron una solución elegante: en lugar de aprender una función $H(x)$ directamente, la red aprende el \textbf{residuo} $F(x) = H(x) - x$.

\begin{equation}
\text{Salida} = F(x) + x \quad \text{(Skip Connection)}
\end{equation}

\begin{figure}[H]
\centering
\begin{tikzpicture}[
    block/.style={rectangle, draw, fill=blue!20, minimum width=2cm, minimum height=0.8cm},
    arrow/.style={->, thick}
]
    % Input
    \node (input) at (0, 0) {$x$};

    % Blocks
    \node[block] (conv1) at (3, 0) {Conv + BN + ReLU};
    \node[block] (conv2) at (3, -1.5) {Conv + BN};

    % Addition
    \node[circle, draw, fill=green!20] (add) at (6, -0.75) {$+$};

    % Output
    \node (output) at (8, -0.75) {ReLU};

    % Arrows
    \draw[arrow] (input) -- (conv1);
    \draw[arrow] (conv1) -- (conv2);
    \draw[arrow] (conv2) -- (add);
    \draw[arrow] (add) -- (output);

    % Skip connection
    \draw[arrow, dashed, red] (input) -- (0, -2.5) -- (6, -2.5) -- (add);

    % Labels
    \node[red] at (3, -2.8) {\small Skip Connection (identidad)};
    \node at (3, -0.9) {\small $F(x)$};

\end{tikzpicture}
\caption{Bloque Residual de ResNet. La conexión skip permite que el gradiente fluya directamente, facilitando el entrenamiento de redes profundas.}
\end{figure}

\subsubsection{¿Por qué ResNet-18?}

La familia ResNet incluye variantes de diferentes profundidades: ResNet-18, 34, 50, 101, 152. Elegimos \textbf{ResNet-18} por las siguientes razones:

\begin{table}[H]
\centering
\caption{Comparación de variantes ResNet}
\label{tab:resnet_variants}
\begin{tabular}{@{}lccl@{}}
\toprule
\textbf{Modelo} & \textbf{Capas} & \textbf{Parámetros} & \textbf{Consideración} \\
\midrule
ResNet-18 & 18 & 11.7M & \cellcolor{green!20}Ideal para datasets pequeños \\
ResNet-34 & 34 & 21.8M & Buen balance \\
ResNet-50 & 50 & 25.6M & Requiere más datos \\
ResNet-101 & 101 & 44.5M & Requiere muchos datos \\
ResNet-152 & 152 & 60.2M & Solo para datasets masivos \\
\bottomrule
\end{tabular}
\end{table}

\begin{tipbox}[Regla Práctica]
Con datasets pequeños ($<$10.000 imágenes), modelos más pequeños como ResNet-18 suelen generalizar mejor. Modelos más profundos tienen mayor capacidad de memorizar el conjunto de entrenamiento (overfitting).
\end{tipbox}

\subsection{DenseNet: Redes Densamente Conectadas}

\subsubsection{Filosofía: Reutilizar Características}

Mientras ResNet suma características mediante skip connections, \textbf{DenseNet} (Huang et al., 2017) lleva esta idea más lejos: cada capa recibe las características de \textbf{todas} las capas anteriores.

\begin{equation}
x_l = H_l([x_0, x_1, ..., x_{l-1}])
\end{equation}

Donde $[...]$ denota concatenación de mapas de características.

\begin{figure}[H]
\centering
\begin{tikzpicture}[
    block/.style={rectangle, draw, fill=orange!30, minimum width=1cm, minimum height=0.6cm},
    arrow/.style={->, thick}
]
    % Blocks
    \node[block] (l0) at (0, 0) {$L_0$};
    \node[block] (l1) at (2, 0) {$L_1$};
    \node[block] (l2) at (4, 0) {$L_2$};
    \node[block] (l3) at (6, 0) {$L_3$};

    % Dense connections
    \draw[arrow] (l0) -- (l1);
    \draw[arrow] (l1) -- (l2);
    \draw[arrow] (l2) -- (l3);

    \draw[arrow, blue] (l0) to[bend left=40] (l2);
    \draw[arrow, blue] (l0) to[bend left=50] (l3);
    \draw[arrow, blue] (l1) to[bend left=40] (l3);

\end{tikzpicture}
\caption{Conectividad densa: cada capa recibe información de todas las capas anteriores.}
\end{figure}

\subsubsection{Ventajas de DenseNet}

\begin{enumerate}
    \item \textbf{Reutilización de características}: Las características de capas tempranas (bordes, texturas) están disponibles directamente en capas tardías
    \item \textbf{Flujo de gradiente mejorado}: Múltiples caminos para el gradiente durante backpropagation
    \item \textbf{Menos parámetros}: Gracias a la reutilización, cada capa puede ser más estrecha
    \item \textbf{Excelente en imágenes médicas}: Publicaciones demuestran su eficacia en radiografías, dermatoscopias, etc.
\end{enumerate}

\subsubsection{¿Por qué DenseNet-121?}

\begin{itemize}
    \item \textbf{Balance tamaño/rendimiento}: 121 capas con ``solo'' 8M de parámetros
    \item \textbf{Probado en contexto médico}: CheXNet (detección de neumonía en rayos X) se basa en DenseNet-121
    \item \textbf{Disponible preentrenado}: En ImageNet, lo que facilita Transfer Learning
\end{itemize}

\subsection{YOLO para Clasificación}

\subsubsection{De Detección a Clasificación}

YOLO (\textit{You Only Look Once}) es famoso como detector de objetos, pero Ultralytics ofrece variantes específicas para \textbf{clasificación de imágenes}. Estas versiones (YOLO11x-cls, YOLO26x-cls) mantienen el backbone convolucional potente de YOLO pero reemplazan las cabezas de detección por capas de clasificación.

\subsubsection{¿Por qué incluir YOLO?}

\begin{enumerate}
    \item \textbf{Arquitectura diferente}: Aporta diversidad al ensemble (no es ResNet ni DenseNet)
    \item \textbf{Pipeline optimizado}: Ultralytics incluye data augmentation avanzado, entrenamiento eficiente y early stopping integrado
    \item \textbf{Estado del arte}: Las versiones recientes (YOLO11, YOLO26) incorporan las últimas técnicas
\end{enumerate}

\subsubsection{YOLO11x vs YOLO26x}

\begin{table}[H]
\centering
\caption{Comparación de modelos YOLO para clasificación}
\begin{tabular}{@{}lccl@{}}
\toprule
\textbf{Modelo} & \textbf{Parámetros} & \textbf{F1 Test} & \textbf{Observación} \\
\midrule
YOLO11x-cls & $\approx$57M & 0.57 & Más rápido \\
YOLO26x-cls & $\approx$90M & \textbf{0.60} & +3\% rendimiento \\
\bottomrule
\end{tabular}
\end{table}

Optamos por \textbf{YOLO26x} en el ensemble final dado su mejor rendimiento, a pesar de su mayor coste computacional.

\subsection{Transfer Learning: Aprovechando Conocimiento Previo}

\subsubsection{El Problema de los Datos Limitados}

Entrenar una CNN desde cero requiere típicamente cientos de miles de imágenes. Con solo 500 imágenes de entrenamiento, la red no puede aprender representaciones robustas y memorizará el conjunto de entrenamiento (overfitting severo).

\subsubsection{La Solución: Pesos Preentrenados}

\textbf{Transfer Learning} consiste en utilizar una red ya entrenada en un dataset grande (típicamente ImageNet con 1.4 millones de imágenes) y adaptarla a nuestra tarea específica.

\begin{conceptbox}[¿Por qué funciona Transfer Learning?]
Las capas iniciales de una CNN aprenden características \textbf{universales}:
\begin{itemize}
    \item \textbf{Capa 1}: Detectores de bordes horizontales, verticales, diagonales
    \item \textbf{Capa 2-3}: Texturas, patrones repetitivos
    \item \textbf{Capas medias}: Partes de objetos (ojos, ruedas, etc.)
    \item \textbf{Capas finales}: Conceptos específicos del dataset original
\end{itemize}
Los bordes y texturas son útiles para \textbf{cualquier} tarea de visión, incluida la dermatoscopia.
\end{conceptbox}

\subsubsection{Estrategia de Fine-Tuning}

Nuestra estrategia de fine-tuning tiene dos componentes:

\begin{enumerate}
    \item \textbf{Congelar capas tempranas}: Mantener fijos los pesos de las primeras capas (características universales)
    \item \textbf{Descongelar capas finales}: Permitir que las últimas capas se adapten a la dermatoscopia
    \item \textbf{Reemplazar el clasificador}: Cambiar la capa final de 1000 clases (ImageNet) a 6 clases (nuestro problema)
\end{enumerate}

\begin{lstlisting}[caption={Estrategia de Fine-Tuning en ResNet-18}]
# 1. Congelar todo
for param in model.parameters():
    param.requires_grad = False

# 2. Descongelar solo layer4 (ultimo bloque)
for param in model.layer4.parameters():
    param.requires_grad = True

# 3. Nueva cabeza para 6 clases
model.fc = nn.Sequential(
    nn.Dropout(0.5),
    nn.Linear(512, 6)
)
\end{lstlisting}

\subsubsection{Learning Rates Diferenciales}

Una técnica avanzada que mejora el fine-tuning es usar \textbf{learning rates diferentes} para distintas partes de la red:

\begin{itemize}
    \item \textbf{Capas preentrenadas (layer4)}: Learning rate bajo ($10^{-4}$) para ajustes finos sin destruir el conocimiento previo
    \item \textbf{Nueva cabeza (fc)}: Learning rate alto ($10^{-3}$) porque debe aprender desde cero
\end{itemize}

\begin{lstlisting}[caption={Optimizador con learning rates diferenciales}]
optimizer = optim.AdamW([
    {'params': model.layer4.parameters(), 'lr': 1e-4},  # Ajuste fino
    {'params': model.fc.parameters(), 'lr': 1e-3}       # Aprendizaje rapido
], weight_decay=0.01)
\end{lstlisting}

\subsection{Data Augmentation: Creando Variabilidad Artificial}

\subsubsection{¿Qué es Data Augmentation?}

El \textbf{Data Augmentation} consiste en aplicar transformaciones aleatorias a las imágenes durante el entrenamiento para crear ``nuevas'' muestras artificiales. Esto ayuda al modelo a ser invariante a estas transformaciones.

\subsubsection{Transformaciones Aplicadas}

\begin{table}[H]
\centering
\caption{Pipeline de Data Augmentation utilizado}
\begin{tabular}{@{}lp{8cm}@{}}
\toprule
\textbf{Transformación} & \textbf{Justificación} \\
\midrule
Resize a 224×224 & Tamaño estándar para Transfer Learning \\
Random Horizontal Flip (50\%) & Una lesión puede aparecer en cualquier orientación \\
Random Vertical Flip (50\%) & Ídem (particularmente útil en dermatoscopia) \\
HSV Jitter & Simula variaciones de iluminación y color de piel \\
Random Resized Crop (0.5-1.0) & Diferentes escalas y encuadres \\
\bottomrule
\end{tabular}
\end{table}

\subsubsection{¿Por qué usar el Pipeline de Ultralytics?}

Utilizamos \texttt{classify\_augmentations} de Ultralytics en lugar de un pipeline manual de torchvision por varias razones:

\begin{enumerate}
    \item \textbf{Consistencia con YOLO}: El mismo preprocesamiento para todos los modelos del ensemble
    \item \textbf{Optimización}: Pipeline altamente optimizado y probado
    \item \textbf{HSV Jitter}: Transformación en espacio HSV que es más perceptualmente relevante que RGB
\end{enumerate}

\begin{lstlisting}[caption={Pipeline de Data Augmentation}]
from ultralytics.data.augment import classify_augmentations

train_pipeline = classify_augmentations(
    size=224,
    mean=[0.485, 0.456, 0.406],  # Media de ImageNet
    std=[0.229, 0.224, 0.225],    # Desv. estandar de ImageNet
    scale=(0.5, 1.0),
    hflip=0.5,
    vflip=0.5,
    hsv_h=0.015, hsv_s=0.7, hsv_v=0.4
)
\end{lstlisting}

\subsection{Métricas de Evaluación}

\subsubsection{¿Por qué NO usar solo Accuracy?}

El \textbf{Accuracy} (precisión global) es la métrica más intuitiva:

\begin{equation}
\text{Accuracy} = \frac{\text{Predicciones Correctas}}{\text{Total de Predicciones}}
\end{equation}

Sin embargo, en datasets desbalanceados el accuracy es \textbf{engañoso}. Un modelo que siempre prediga la clase mayoritaria tendría alto accuracy pero sería inútil clínicamente.

\begin{warningbox}[Ejemplo del Peligro del Accuracy]
Si el 60\% de las imágenes fueran de clase NEV, un modelo que siempre predijera ``NEV'' tendría 60\% de accuracy, pero fallaría en detectar \textbf{todos} los melanomas. Esto sería inaceptable médicamente.
\end{warningbox}

\subsubsection{Precision, Recall y F1-Score}

Para cada clase $c$, definimos:

\begin{align}
\text{Precision}_c &= \frac{TP_c}{TP_c + FP_c} = \frac{\text{Verdaderos Positivos}}{\text{Predichos como } c} \\[0.3cm]
\text{Recall}_c &= \frac{TP_c}{TP_c + FN_c} = \frac{\text{Verdaderos Positivos}}{\text{Realmente son } c} \\[0.3cm]
\text{F1}_c &= 2 \cdot \frac{\text{Precision}_c \cdot \text{Recall}_c}{\text{Precision}_c + \text{Recall}_c}
\end{align}

\begin{conceptbox}[Interpretación Clínica]
\begin{itemize}
    \item \textbf{Precision}: ``De todos los que diagnostiqué como melanoma, ¿cuántos realmente lo eran?'' (Evita falsos positivos / alarmas innecesarias)
    \item \textbf{Recall (Sensibilidad)}: ``De todos los melanomas reales, ¿cuántos detecté?'' (Evita falsos negativos / casos perdidos)
    \item \textbf{F1-Score}: Media armónica, balancea ambas métricas
\end{itemize}
\end{conceptbox}

\subsubsection{F1-Score Macro: Nuestra Métrica Principal}

El \textbf{F1-Score Macro} es el promedio simple del F1 de cada clase:

\begin{equation}
\text{F1-Macro} = \frac{1}{C} \sum_{c=1}^{C} \text{F1}_c
\end{equation}

Elegimos F1-Macro porque:
\begin{itemize}
    \item \textbf{Da igual importancia a todas las clases}, independientemente de su frecuencia
    \item \textbf{Penaliza fallar en clases minoritarias} como MEL
    \item \textbf{Es la métrica estándar} para clasificación multiclase desbalanceada
\end{itemize}

\subsubsection{Cohen's Kappa}

El \textbf{Coeficiente Kappa de Cohen} mide el acuerdo entre predicciones y etiquetas reales, corrigiendo por el acuerdo que ocurriría por azar:

\begin{equation}
\kappa = \frac{p_o - p_e}{1 - p_e}
\end{equation}

Donde $p_o$ es la proporción de acuerdos observados y $p_e$ es la proporción de acuerdos esperados por azar.

\begin{table}[H]
\centering
\caption{Interpretación del Coeficiente Kappa}
\begin{tabular}{@{}ll@{}}
\toprule
\textbf{Valor de $\kappa$} & \textbf{Interpretación} \\
\midrule
$< 0$ & Peor que el azar \\
$0.00 - 0.20$ & Acuerdo insignificante \\
$0.21 - 0.40$ & Acuerdo regular \\
$0.41 - 0.60$ & Acuerdo moderado \\
$0.61 - 0.80$ & Acuerdo sustancial \\
$0.81 - 1.00$ & Acuerdo casi perfecto \\
\bottomrule
\end{tabular}
\end{table}

\subsection{Técnicas de Regularización}

La regularización es crucial para evitar el \textbf{overfitting}, especialmente con datasets pequeños.

\subsubsection{Dropout}

El \textbf{Dropout} (Srivastava et al., 2014) consiste en ``apagar'' aleatoriamente un porcentaje de neuronas durante el entrenamiento.

\begin{equation}
\hat{y}_i = \begin{cases}
0 & \text{con probabilidad } p \\
\frac{y_i}{1-p} & \text{con probabilidad } 1-p
\end{cases}
\end{equation}

\begin{itemize}
    \item \textbf{Efecto}: Impide que la red dependa excesivamente de neuronas específicas
    \item \textbf{Interpretación}: Entrena implícitamente un ensemble de sub-redes
    \item \textbf{Valor usado}: $p = 0.5$ (50\% de neuronas apagadas)
\end{itemize}

\subsubsection{Weight Decay (Regularización L2)}

El \textbf{Weight Decay} añade un término de penalización a la función de pérdida que castiga pesos grandes:

\begin{equation}
\mathcal{L}_{\text{total}} = \mathcal{L}_{\text{CE}} + \lambda \sum_i w_i^2
\end{equation}

Donde $\lambda$ es el coeficiente de weight decay (usamos $\lambda = 0.01$).

\begin{itemize}
    \item \textbf{Efecto}: Favorece pesos pequeños, produciendo modelos más ``suaves''
    \item \textbf{Beneficio}: Reduce la capacidad de memorizar el conjunto de entrenamiento
\end{itemize}

\subsection{Optimizador: AdamW}

Utilizamos \textbf{AdamW} (Loshchilov \& Hutter, 2017), una variante de Adam que implementa correctamente el weight decay.

\begin{conceptbox}[¿Por qué AdamW y no SGD o Adam?]
\begin{itemize}
    \item \textbf{vs SGD}: AdamW adapta el learning rate por parámetro, convergiendo más rápido
    \item \textbf{vs Adam clásico}: Adam implementaba weight decay incorrectamente (lo aplicaba al gradiente adaptado). AdamW lo aplica directamente a los pesos, lo cual es teóricamente correcto
\end{itemize}
\end{conceptbox}

\subsection{Función de Pérdida: Cross-Entropy Ponderada}

\subsubsection{Cross-Entropy Estándar}

Para clasificación multiclase, la pérdida estándar es \textbf{Cross-Entropy}:

\begin{equation}
\mathcal{L}_{\text{CE}} = -\sum_{c=1}^{C} y_c \log(\hat{y}_c)
\end{equation}

Donde $y_c$ es la etiqueta real (one-hot) y $\hat{y}_c$ es la probabilidad predicha para la clase $c$.

\subsubsection{Ponderación por Clase}

Para datasets desbalanceados, ponderamos la pérdida de cada clase inversamente proporcional a su frecuencia:

\begin{equation}
w_c = \frac{1}{n_c} \cdot \frac{N}{C}
\end{equation}

Donde $n_c$ es el número de muestras de la clase $c$, $N$ es el total de muestras, y $C$ es el número de clases.

\begin{lstlisting}[caption={Cálculo de pesos de clase}]
class_counts = np.bincount(labels)
weights = 1.0 / class_counts
weights = weights / weights.sum() * len(class_counts)
# Resultado: [0.75, 0.75, 2.25, 0.75, 0.75, 0.75]
# MEL (indice 2) tiene peso 3x mayor
\end{lstlisting}

\subsection{Ensemble: Combinando Múltiples Modelos}

\subsubsection{¿Por qué un Ensemble?}

Un \textbf{ensemble} combina las predicciones de múltiples modelos. Esto funciona porque diferentes arquitecturas cometen \textbf{errores diferentes}. Al promediar, los errores individuales tienden a cancelarse.

\subsubsection{Estrategia: Promedio Ponderado de Probabilidades}

Nuestra estrategia de fusión es un \textbf{promedio ponderado} (soft voting):

\begin{equation}
P_{\text{ensemble}}(c) = \sum_{m=1}^{M} \alpha_m \cdot P_m(c) \quad \text{donde } \sum_{m} \alpha_m = 1
\end{equation}

Los pesos $\alpha_m$ se asignan según el rendimiento individual de cada modelo:

\begin{table}[H]
\centering
\caption{Pesos del ensemble final}
\begin{tabular}{@{}lcc@{}}
\toprule
\textbf{Modelo} & \textbf{F1 Individual} & \textbf{Peso en Ensemble} \\
\midrule
YOLO26x & 0.60 & 60\% \\
DenseNet-121 & 0.52 & 20\% \\
ResNet-18 & 0.56 & 20\% \\
\bottomrule
\end{tabular}
\end{table}

\begin{tipbox}[Principio del Ensemble]
El modelo con mejor rendimiento individual (YOLO26x) recibe el mayor peso. Los otros modelos aportan ``diversidad'' que puede corregir errores específicos de YOLO.
\end{tipbox}

% ============================================
% 3. METODOLOGÍA
% ============================================
\section{Metodología}

\subsection{Pipeline General}

El pipeline de experimentación sigue las siguientes fases:

\begin{enumerate}
    \item \textbf{Análisis Exploratorio}: Distribución de clases, tamaños de imagen
    \item \textbf{Preprocesamiento}: Data augmentation, normalización ImageNet
    \item \textbf{Entrenamiento Individual}: ResNet-18, DenseNet-121, YOLO26x
    \item \textbf{Evaluación Individual}: Métricas por modelo en el conjunto de test
    \item \textbf{Construcción del Ensemble}: Combinación ponderada
    \item \textbf{Evaluación Final}: Métricas del ensemble
\end{enumerate}

\subsection{Configuración del Entrenamiento}

\begin{table}[H]
\centering
\caption{Hiperparámetros de entrenamiento}
\begin{tabular}{@{}ll@{}}
\toprule
\textbf{Parámetro} & \textbf{Valor} \\
\midrule
Tamaño de imagen & 224 × 224 \\
Batch size & 64 \\
Épocas máximas & 25 \\
Learning rate (capas preentrenadas) & $10^{-4}$ \\
Learning rate (cabeza nueva) & $10^{-3}$ \\
Optimizador & AdamW \\
Weight decay & 0.01 \\
Dropout & 0.5 \\
Early stopping (YOLO) & patience = 5 \\
Criterio de guardado & Mejor F1-Macro en validación \\
\bottomrule
\end{tabular}
\end{table}

\subsection{Estrategia de Fine-Tuning por Arquitectura}

\begin{table}[H]
\centering
\caption{Capas descongeladas por arquitectura}
\begin{tabular}{@{}lll@{}}
\toprule
\textbf{Modelo} & \textbf{Capas Descongeladas} & \textbf{Razón} \\
\midrule
ResNet-18 & layer4 + fc & Último bloque residual \\
DenseNet-121 & denseblock4 + norm5 + classifier & Último bloque denso \\
YOLO26x & Entrenamiento completo & Pipeline Ultralytics \\
\bottomrule
\end{tabular}
\end{table}

\subsection{Normalización}

Todas las imágenes se normalizan con la media y desviación estándar de ImageNet:

\begin{align}
\mu &= [0.485, 0.456, 0.406] \\
\sigma &= [0.229, 0.224, 0.225]
\end{align}

Esto es necesario porque los modelos preentrenados esperan entradas normalizadas de esta forma.

% ============================================
% 4. RESULTADOS
% ============================================
\section{Resultados Experimentales}

\subsection{Análisis del Dataset}

\begin{figure}[H]
\centering
% Placeholder para imagen
\fbox{\parbox{0.8\textwidth}{\centering\vspace{2cm}\textbf{[FIGURA: Distribución de clases por conjunto (train/val/test)]}\vspace{2cm}}}
\caption{Distribución de clases en cada subconjunto. Nótese el marcado desbalanceo en MEL.}
\label{fig:class_distribution}
\end{figure}

\subsection{Resultados por Modelo Individual}

\subsubsection{ResNet-18}

\begin{table}[H]
\centering
\caption{Resultados de ResNet-18 en el conjunto de test}
\begin{tabular}{@{}lcccc@{}}
\toprule
\textbf{Clase} & \textbf{Precision} & \textbf{Recall} & \textbf{F1-Score} & \textbf{Support} \\
\midrule
ACK & 0.45 & 0.50 & 0.47 & 30 \\
BCC & 0.50 & 0.47 & 0.48 & 30 \\
\rowcolor{red!10} MEL & 0.75 & 0.55 & 0.63 & 11 \\
NEV & 0.65 & 0.57 & 0.61 & 30 \\
SCC & 0.44 & 0.37 & 0.40 & 30 \\
SEK & 0.63 & 0.83 & 0.72 & 30 \\
\midrule
\textbf{Macro avg} & 0.57 & 0.55 & \textbf{0.55} & 161 \\
\bottomrule
\end{tabular}
\end{table}

\textbf{Mejor F1 en Validación}: 0.5596

\subsubsection{DenseNet-121}

\begin{table}[H]
\centering
\caption{Resultados de DenseNet-121 en el conjunto de test}
\begin{tabular}{@{}lcccc@{}}
\toprule
\textbf{Clase} & \textbf{Precision} & \textbf{Recall} & \textbf{F1-Score} & \textbf{Support} \\
\midrule
ACK & 0.55 & 0.37 & 0.44 & 30 \\
BCC & 0.48 & 0.53 & 0.51 & 30 \\
\rowcolor{red!10} MEL & 0.50 & 0.45 & 0.48 & 11 \\
NEV & 0.56 & 0.50 & 0.53 & 30 \\
SCC & 0.39 & 0.47 & 0.42 & 30 \\
SEK & 0.61 & 0.77 & 0.68 & 30 \\
\midrule
\textbf{Macro avg} & 0.52 & 0.51 & \textbf{0.51} & 161 \\
\bottomrule
\end{tabular}
\end{table}

\textbf{Mejor F1 en Validación}: 0.5997

\subsubsection{YOLO26x}

\begin{table}[H]
\centering
\caption{Resultados de YOLO26x en el conjunto de test}
\begin{tabular}{@{}lcccc@{}}
\toprule
\textbf{Clase} & \textbf{Precision} & \textbf{Recall} & \textbf{F1-Score} & \textbf{Support} \\
\midrule
ACK & 0.58 & 0.50 & 0.54 & 30 \\
BCC & 0.58 & 0.47 & 0.52 & 30 \\
\rowcolor{green!20} MEL & \textbf{1.00} & 0.55 & \textbf{0.71} & 11 \\
NEV & 0.67 & 0.60 & 0.63 & 30 \\
SCC & 0.43 & 0.50 & 0.46 & 30 \\
SEK & 0.65 & 0.93 & 0.77 & 30 \\
\midrule
\textbf{Macro avg} & 0.65 & 0.59 & \textbf{0.60} & 161 \\
\bottomrule
\end{tabular}
\end{table}

\begin{tipbox}[Resultado Destacado: MEL]
YOLO26x logra \textbf{100\% de precisión en Melanoma}. Esto significa que cuando el modelo predice melanoma, \textbf{siempre acierta}. No hay falsos positivos, evitando alarmas innecesarias a pacientes.
\end{tipbox}

\subsection{Comparativa de Modelos Individuales}

\begin{table}[H]
\centering
\caption{Resumen comparativo de modelos individuales}
\begin{tabular}{@{}lccccc@{}}
\toprule
\textbf{Modelo} & \textbf{Val F1} & \textbf{Test Acc} & \textbf{Test F1-Macro} & \textbf{MEL F1} & \textbf{Kappa} \\
\midrule
ResNet-18 & 0.5596 & 0.55 & 0.55 & 0.63 & 0.46 \\
DenseNet-121 & 0.5997 & 0.52 & 0.51 & 0.48 & 0.42 \\
\rowcolor{green!20} YOLO26x & --- & \textbf{0.60} & \textbf{0.60} & \textbf{0.71} & \textbf{0.52} \\
\bottomrule
\end{tabular}
\end{table}

\textbf{Observaciones}:
\begin{itemize}
    \item YOLO26x es el mejor modelo individual en todas las métricas
    \item Existe discrepancia entre F1 de validación y test en DenseNet (posible overfitting)
    \item Todos los modelos superan el umbral de ``acuerdo moderado'' en Kappa ($\kappa > 0.4$)
\end{itemize}

\subsection{Resultados del Ensemble}

\begin{table}[H]
\centering
\caption{Resultados del Ensemble (ResNet + DenseNet + YOLO26x)}
\begin{tabular}{@{}lcccc@{}}
\toprule
\textbf{Clase} & \textbf{Precision} & \textbf{Recall} & \textbf{F1-Score} & \textbf{Support} \\
\midrule
ACK & 0.67 & 0.47 & 0.55 & 30 \\
BCC & 0.54 & 0.47 & 0.50 & 30 \\
\rowcolor{green!20} MEL & \textbf{1.00} & 0.55 & \textbf{0.71} & 11 \\
NEV & 0.72 & 0.70 & 0.71 & 30 \\
SCC & 0.45 & 0.57 & 0.50 & 30 \\
SEK & 0.71 & 0.97 & 0.82 & 30 \\
\midrule
\textbf{Macro avg} & \textbf{0.68} & 0.62 & \textbf{0.63} & 161 \\
\bottomrule
\end{tabular}
\end{table}

\begin{table}[H]
\centering
\caption{Métricas globales del Ensemble}
\begin{tabular}{@{}ll@{}}
\toprule
\textbf{Métrica} & \textbf{Valor} \\
\midrule
Accuracy Global & 0.6273 (62.7\%) \\
F1-Score Macro & 0.63 \\
Precision Macro & 0.68 \\
Recall Macro & 0.62 \\
Cohen's Kappa & 0.55 \\
\bottomrule
\end{tabular}
\end{table}

\subsection{Análisis de la Matriz de Confusión}

\begin{figure}[H]
\centering
% Placeholder para imagen
\fbox{\parbox{0.7\textwidth}{\centering\vspace{5cm}\textbf{[FIGURA: Matriz de Confusión del Ensemble]}\vspace{5cm}}}
\caption{Matriz de confusión del ensemble en el conjunto de test. Los valores en la diagonal representan predicciones correctas.}
\label{fig:confusion_matrix}
\end{figure}

\textbf{Observaciones de la matriz de confusión}:
\begin{itemize}
    \item \textbf{SEK} es la clase mejor clasificada (29/30 correctos, F1=0.82)
    \item \textbf{MEL} mantiene 100\% de precisión (0 falsos positivos)
    \item Las confusiones principales ocurren entre:
    \begin{itemize}
        \item ACK $\leftrightarrow$ SCC (ambas queratosis, visualmente similares)
        \item BCC $\leftrightarrow$ NEV (carcinoma vs lunar)
    \end{itemize}
\end{itemize}

\subsection{Curvas de Entrenamiento}

\begin{figure}[H]
\centering
% Placeholder para imagen
\fbox{\parbox{0.9\textwidth}{\centering\vspace{4cm}\textbf{[FIGURA: Evolución de Loss y F1 durante el entrenamiento de ResNet-18]}\vspace{4cm}}}
\caption{Curvas de entrenamiento de ResNet-18. Izquierda: Train/Val Loss. Derecha: Val F1-Score.}
\label{fig:training_curves}
\end{figure}

\textbf{Observaciones}:
\begin{itemize}
    \item El train loss desciende consistentemente
    \item El val loss empieza a subir alrededor de la época 15 (inicio de overfitting)
    \item El mejor F1 de validación se alcanza antes del overfitting severo
    \item El mecanismo de ``guardar el mejor modelo'' previene usar pesos sobreajustados
\end{itemize}

% ============================================
% 5. DISCUSIÓN
% ============================================
\section{Discusión}

\subsection{Análisis del Rendimiento}

\subsubsection{¿Por qué YOLO26x es el mejor modelo?}

YOLO26x supera a las arquitecturas clásicas por varias razones:

\begin{enumerate}
    \item \textbf{Arquitectura más moderna}: Incorpora técnicas de 2024-2025 (attention mechanisms, depthwise separable convolutions, etc.)
    \item \textbf{Pipeline de entrenamiento optimizado}: Early stopping, data augmentation avanzado, warm-up de learning rate
    \item \textbf{Mayor capacidad}: Con $\approx$90M de parámetros, puede capturar patrones más complejos
\end{enumerate}

\subsubsection{¿Aporta el Ensemble Valor?}

\begin{table}[H]
\centering
\caption{Comparación: Mejor modelo individual vs Ensemble}
\begin{tabular}{@{}lccc@{}}
\toprule
\textbf{Sistema} & \textbf{F1-Macro} & \textbf{Precision} & \textbf{MEL F1} \\
\midrule
YOLO26x (solo) & 0.60 & 0.65 & 0.71 \\
\rowcolor{green!20} Ensemble & \textbf{0.63} & \textbf{0.68} & 0.71 \\
\midrule
\textbf{Mejora} & \textbf{+3\%} & \textbf{+3\%} & = \\
\bottomrule
\end{tabular}
\end{table}

El ensemble aporta una \textbf{mejora de 3 puntos porcentuales} en F1-Macro y Precision, manteniendo el excelente rendimiento en MEL. Esta mejora, aunque modesta, es valiosa en contexto médico donde cada punto de precisión puede significar diagnósticos correctos adicionales.

\subsection{Limitaciones del Estudio}

\begin{enumerate}
    \item \textbf{Tamaño del dataset}: Con solo 500 imágenes de entrenamiento, los modelos tienen capacidad limitada de generalización

    \item \textbf{Desbalanceo extremo en MEL}: Solo $\approx$35 imágenes de melanoma para entrenamiento dificultan el aprendizaje de características discriminativas

    \item \textbf{Overfitting}: A pesar de las técnicas de regularización, se observa divergencia train/val loss

    \item \textbf{Validación externa}: Los resultados son específicos de este dataset; se requeriría validación en datasets externos para confirmar generalización
\end{enumerate}

\subsection{Comparación con el Estado del Arte}

Los modelos de clasificación de lesiones cutáneas en competiciones ISIC alcanzan F1-Macro superiores al 85\% en datasets más grandes. Nuestro 63\% refleja principalmente la limitación del tamaño del dataset, no deficiencias metodológicas.

% ============================================
% 6. CONCLUSIONES
% ============================================
\section{Conclusiones}

\subsection{Resumen de Logros}

\begin{enumerate}
    \item Se implementaron y compararon tres arquitecturas de Deep Learning (ResNet-18, DenseNet-121, YOLO26x) para clasificación de lesiones cutáneas

    \item Se aplicó exitosamente \textbf{Transfer Learning}, fundamental dado el tamaño limitado del dataset

    \item Se construyó un \textbf{ensemble} que mejora el rendimiento del mejor modelo individual en 3 puntos porcentuales

    \item Se logró \textbf{100\% de precisión en la detección de Melanoma}, la clase clínicamente más crítica

    \item Se documentaron y justificaron todas las decisiones técnicas desde una perspectiva tanto teórica como práctica
\end{enumerate}

\subsection{Métricas Finales}

\begin{table}[H]
\centering
\caption{Resultados finales del sistema}
\begin{tabular}{@{}ll@{}}
\toprule
\textbf{Métrica} & \textbf{Valor} \\
\midrule
\textbf{Accuracy} & 62.7\% \\
\textbf{F1-Score Macro} & 0.63 \\
\textbf{Precision Macro} & 0.68 \\
\textbf{Recall Macro} & 0.62 \\
\textbf{Cohen's Kappa} & 0.55 (Acuerdo Moderado) \\
\textbf{MEL F1-Score} & 0.71 \\
\textbf{MEL Precision} & 1.00 \\
\bottomrule
\end{tabular}
\end{table}

\subsection{Lecciones Aprendidas}

\begin{itemize}
    \item \textbf{El tamaño del dataset es crítico}: La principal limitación no fue la arquitectura sino la cantidad de datos

    \item \textbf{Transfer Learning es esencial}: Sin pesos preentrenados, el overfitting sería mucho peor

    \item \textbf{Las métricas correctas importan}: F1-Macro reveló problemas que Accuracy ocultaba

    \item \textbf{Los ensembles aportan robustez}: Incluso modelos individualmente ``peores'' contribuyen positivamente

    \item \textbf{YOLO no es solo para detección}: Las versiones de clasificación son competitivas con arquitecturas clásicas
\end{itemize}

\subsection{Trabajo Futuro}

Este trabajo establece la línea base para la parte de Deep Learning. En las siguientes fases del proyecto integrado se:

\begin{enumerate}
    \item Integrarán los \textbf{datos tabulares} (metadatos del paciente) mediante modelos de Machine Learning clásico

    \item Implementará la \textbf{fusión} de ambas modalidades mediante:
    \begin{itemize}
        \item Fusión tardía (stacking/ensemble)
        \item Fusión temprana (embeddings de imagen + features tabulares)
    \end{itemize}

    \item Evaluará si la combinación \textbf{imagen + metadatos} supera a cada modalidad por separado
\end{enumerate}

% ============================================
% REFERENCIAS
% ============================================
\section*{Referencias}

\begin{enumerate}
    \item He, K., Zhang, X., Ren, S., \& Sun, J. (2016). Deep residual learning for image recognition. In \textit{CVPR}.

    \item Huang, G., Liu, Z., Van Der Maaten, L., \& Weinberger, K. Q. (2017). Densely connected convolutional networks. In \textit{CVPR}.

    \item Loshchilov, I., \& Hutter, F. (2017). Decoupled weight decay regularization. \textit{arXiv preprint arXiv:1711.05101}.

    \item Srivastava, N., Hinton, G., Krizhevsky, A., Sutskever, I., \& Salakhutdinov, R. (2014). Dropout: a simple way to prevent neural networks from overfitting. \textit{JMLR}.

    \item Ultralytics (2024). YOLO Documentation. \url{https://docs.ultralytics.com}

    \item ISIC Archive. International Skin Imaging Collaboration. \url{https://www.isic-archive.com}
\end{enumerate}

% ============================================
% APÉNDICE
% ============================================
\appendix
\section{Código de Entrenamiento}

A continuación se presenta el código completo utilizado para entrenar los modelos.

\subsection{Función de Entrenamiento}

\begin{lstlisting}[caption={Función principal de entrenamiento}]
def entrenar_modelo(model, train_loader, val_loader, criterion,
                    optimizer, device, epochs=25, nombre_archivo="modelo.pth"):

    historial = {'train_loss': [], 'val_loss': [], 'val_f1': []}
    best_f1 = 0.0
    best_model_wts = copy.deepcopy(model.state_dict())

    for epoch in range(epochs):
        # FASE ENTRENAMIENTO
        model.train()
        running_loss = 0.0

        for inputs, labels in train_loader:
            inputs, labels = inputs.to(device), labels.to(device)

            optimizer.zero_grad()
            outputs = model(inputs)
            loss = criterion(outputs, labels)
            loss.backward()
            optimizer.step()

            running_loss += loss.item() * inputs.size(0)

        epoch_train_loss = running_loss / len(train_loader.dataset)

        # FASE VALIDACION
        model.eval()
        val_loss = 0.0
        y_true, y_pred = [], []

        with torch.no_grad():
            for inputs, labels in val_loader:
                inputs, labels = inputs.to(device), labels.to(device)
                outputs = model(inputs)
                loss = criterion(outputs, labels)
                val_loss += loss.item() * inputs.size(0)

                _, preds = torch.max(outputs, 1)
                y_true.extend(labels.cpu().numpy())
                y_pred.extend(preds.cpu().numpy())

        epoch_val_loss = val_loss / len(val_loader.dataset)
        epoch_val_f1 = f1_score(y_true, y_pred, average='macro')

        # GUARDAR MEJOR MODELO
        if epoch_val_f1 > best_f1:
            best_f1 = epoch_val_f1
            best_model_wts = copy.deepcopy(model.state_dict())
            torch.save(model.state_dict(), nombre_archivo)

    model.load_state_dict(best_model_wts)
    return model, historial
\end{lstlisting}

\subsection{Configuración de ResNet-18}

\begin{lstlisting}[caption={Configuración de ResNet-18 para fine-tuning}]
# Cargar modelo preentrenado
model = models.resnet18(weights=models.ResNet18_Weights.DEFAULT)

# Congelar capas
for param in model.parameters():
    param.requires_grad = False

# Descongelar layer4
for param in model.layer4.parameters():
    param.requires_grad = True

# Nueva cabeza
num_ftrs = model.fc.in_features
model.fc = nn.Sequential(
    nn.Dropout(0.5),
    nn.Linear(num_ftrs, 6)  # 6 clases
)

# Optimizador con learning rates diferenciales
optimizer = optim.AdamW([
    {'params': model.layer4.parameters(), 'lr': 1e-4},
    {'params': model.fc.parameters(), 'lr': 1e-3}
], weight_decay=0.01)
\end{lstlisting}

\subsection{Código del Ensemble}

\begin{lstlisting}[caption={Implementación del ensemble}]
# Pesos del ensemble
WEIGHTS = {'resnet': 0.20, 'densenet': 0.20, 'yolo': 0.60}

for img_path in test_images:
    # Predicciones PyTorch
    img_tensor = transform(Image.open(img_path)).unsqueeze(0).to(device)

    with torch.no_grad():
        probs_resnet = F.softmax(model_resnet(img_tensor), dim=1)
        probs_densenet = F.softmax(model_densenet(img_tensor), dim=1)

    # Prediccion YOLO
    result_yolo = model_yolo(img_path, verbose=False)[0]
    probs_yolo = result_yolo.probs.data

    # Fusion ponderada
    probs_final = (
        WEIGHTS['resnet'] * probs_resnet +
        WEIGHTS['densenet'] * probs_densenet +
        WEIGHTS['yolo'] * probs_yolo
    )

    prediction = probs_final.argmax()
\end{lstlisting}

\end{document}
